\documentclass[11pt,a4paper]{article}
\usepackage[utf8]{inputenc}
\usepackage[T1]{fontenc}
\usepackage{times}
\usepackage{graphicx}
\usepackage{amsmath,amssymb}
\usepackage{booktabs}
\usepackage{multirow}
\usepackage{xcolor}
\usepackage[margin=2.5cm]{geometry}
\usepackage{hyperref}
\usepackage{cleveref}
\usepackage{natbib}
\bibliographystyle{plainnat}

\newcommand{\cmark}{\checkmark}

\title{Causal Style Decomposition Reveals Systematic Style Bias\\in LLM-as-Judge That Likert Scoring Misses}

\author{Anonymous}
\date{}

\begin{document}
\maketitle

% ============================================================
% ABSTRACT
% ============================================================
\begin{abstract}
LLM-as-judge has become critical infrastructure for NLP evaluation and alignment training, yet recent work shows style bias dominates other judge biases.
Existing analyses treat style as a monolithic dimension or rely on correlational evidence.
We introduce \textbf{Causal Style Decomposition (CSD)}, a black-box, post-hoc framework that generates binary style counterfactuals, verifies content preservation via NLI and BERTScore, and measures causal style bias through position-controlled pairwise comparison.
Across 3 judges (GPT-4o, Claude-Sonnet, Gemini-Flash), 3 benchmarks (LLMBar, AlpacaEval, MT-Bench), and 3 style axes (formality, tone, register), CSD detects significant style premiums (mean win rate 66.3\%, $p<0.05$ in 7--9 of 9 judge$\times$benchmark cells for formality).
We uncover a systematic \emph{measurement gap}: standard Likert pointwise scoring fails to detect any of these biases (0/9 significant), while pairwise comparison succeeds (7/9 significant)---implying that prior Likert-based studies may systematically underestimate style bias severity.
Style bias is axis-dependent: formality (68.4\%) $>$ register (66.3\%) $>$ tone (63.4\%), with lighter models (Gemini-Flash) consistently most susceptible.
A scaled replication ($n=500$, $p<10^{-6}$), original-preference control, and rewriter ablation confirm robustness.
These findings have direct implications for evaluation methodology: style bias audits should use pairwise comparison, and bias must be assessed per axis rather than as a single ``style'' factor.
\end{abstract}

% ============================================================
% 1. INTRODUCTION
% ============================================================
\section{Introduction}

LLM-as-judge---using large language models to evaluate text quality---has become a cornerstone of NLP evaluation and preference-based alignment training.
Systems like AlpacaEval, MT-Bench, and Chatbot Arena rely on LLM judges to rank model outputs, and RLHF pipelines use judge-generated preferences as reward signals.
The reliability of these judgments directly affects which models are deemed superior and which behaviors are reinforced during training.

Recent studies have established that \emph{style bias} is the dominant source of unreliability in LLM judges.
SOS-Bench shows style features explain judge scores with $R=0.999$, far exceeding content features.
Preference Leakage demonstrates that surface-level style cues drive 76--92\% of observed preference signals.
Yet existing approaches to style bias have critical limitations:
(1)~they typically treat style as a single dimension, usually length;
(2)~they rely on correlational analyses that cannot establish causation;
(3)~population-level corrections (e.g., length-controlled AlpacaEval) cannot diagnose per-instance bias;
(4)~no existing method provides black-box, post-hoc, per-axis causal style bias detection.

We introduce \textbf{Causal Style Decomposition (CSD)}, a framework that addresses all four gaps.
CSD generates binary counterfactual pairs along a single style axis (e.g., formal vs.\ casual), verifies content preservation through NLI entailment and BERTScore, and measures the judge's causal response to the style intervention via position-controlled pairwise comparison.
CSD is model-agnostic, requires no access to judge internals, and applies to any style dimension where content-preserving rewriting is feasible.

Our experiments yield three main findings:
\begin{enumerate}
    \item \textbf{C1 (Method Validity):} CSD detects significant causal style bias across 3 judges, 3 benchmarks, and 3 style axes, with formality premium ranging from 60--72\% win rate ($p<0.05$).
    \item \textbf{C2 (Measurement Insight):} Standard Likert pointwise scoring \emph{systematically} fails to detect style biases that pairwise comparison reveals. In 7 of 9 judge$\times$benchmark combinations, Likert yields $p>0.10$ while pairwise yields $p<0.05$---a measurement gap with direct implications for evaluation methodology.
    \item \textbf{C3 (Axis Dependency):} Style bias is not monolithic. Different axes produce different effect sizes (formality $>$ register $>$ tone), and different judges show different sensitivity profiles (Gemini-Flash consistently strongest, Claude-Sonnet weakest).
\end{enumerate}

% ============================================================
% 2. RELATED WORK
% ============================================================
\section{Related Work}

\paragraph{LLM-as-Judge Bias.}
Position bias---the tendency to prefer responses in certain positions---has been extensively studied.
More recent work shifts attention to style bias as the dominant factor.
SOS-Bench~injects quality defects to measure judge sensitivity to style vs.\ content, finding style overwhelmingly dominant.
Preference Leakage~demonstrates that RLHF preference data encode surface-style cues that leak into judge behavior.
CASSE~constructs a 12$\times$12 style feature interaction matrix but requires prompt engineering for each feature, creating causal entanglement issues.

\paragraph{Length Bias Correction.}
Length-controlled AlpacaEval~applies regression-based correction for response length, the most-studied single style dimension.
While effective for length, this approach does not generalize to other style axes and operates at population level rather than per-instance.

\paragraph{Counterfactual Evaluation.}
Counterfactual data augmentation has been used in NLI, sentiment analysis, and fairness auditing.
We adapt this paradigm to LLM judge evaluation: generating style-only counterfactuals and measuring the judge's differential response as a causal effect estimate.

\paragraph{Causal Inference in NLP.}
The use of interventional methods to identify causal effects in NLP systems is a growing trend.
CSD contributes a practical instantiation for the specific problem of LLM judge style bias, combining counterfactual generation with a paired comparison design that yields clean causal estimates.

% ============================================================
% 3. METHOD
% ============================================================
\section{Method: Causal Style Decomposition}

CSD operates as a four-stage pipeline applied post-hoc to any black-box LLM judge.

\subsection{Stage 1: Counterfactual Generation}

Given a response $r$ and a target style axis $a$ with poles $(s^+, s^-)$ (e.g., formal/casual), we generate a counterfactual $r'$ by rewriting $r$ to shift from its current style pole to the opposite, while preserving semantic content.
We use an instruction-tuned LLM (Qwen3-32B) as the rewriter with a carefully designed prompt that specifies: (1)~the target style transformation, (2)~explicit content-preservation constraints, and (3)~examples of acceptable vs.\ unacceptable rewrites.

For each original response $r_i$, we produce both directions: $r_i^{s^+}$ (style $s^+$ version) and $r_i^{s^-}$ (style $s^-$ version).
One of these is the original; the other is the counterfactual rewrite.

\subsection{Stage 2: Content Verification}

Not all rewrites successfully preserve content.
We apply two complementary filters:
\begin{itemize}
    \item \textbf{NLI Entailment}: We compute bidirectional NLI entailment scores between $r$ and $r'$. Pairs with entailment score $< 0.85$ are excluded.
    \item \textbf{BERTScore}: We compute BERTScore F1 between $r$ and $r'$ as a secondary semantic similarity check.
\end{itemize}

Across our experiments, NLI pass rates range from 86--90\% depending on the axis (\Cref{tab:quality}), confirming that content-preserving style rewriting is feasible for surface-level style dimensions.

\subsection{Stage 3: Position-Controlled Pairwise Comparison}

For each verified counterfactual pair $(r_i^{s^+}, r_i^{s^-})$, we present the pair to the judge in both orders (A/B and B/A) and record preferences.
A pair is \emph{consistent} if the judge selects the same style pole regardless of position.
Inconsistent pairs (position-driven) are excluded from the style bias estimate.

This dual-direction protocol controls for position bias, which is known to affect LLM judges.

\subsection{Stage 4: Statistical Analysis}

On the consistent subset, we compute:
\begin{itemize}
    \item \textbf{Pairwise win rate}: the fraction of consistent pairs where style $s^+$ is preferred.
    \item \textbf{Sign test}: exact binomial test against the null hypothesis of 50\% win rate.
    \item \textbf{Wilson confidence intervals}: for uncertainty quantification.
\end{itemize}

For comparison, we also compute \textbf{Likert ATE} (Average Treatment Effect): each response is scored independently on a Likert scale, and we compute the mean score difference between $s^+$ and $s^-$ versions, tested via paired $t$-test.

% ============================================================
% 4. EXPERIMENTAL SETUP
% ============================================================
\section{Experimental Setup}

\paragraph{Judges.}
We evaluate three commercial LLM judges spanning different model families: \textbf{GPT-4o} (OpenAI), \textbf{Claude-Sonnet} (Anthropic), and \textbf{Gemini-2.0-Flash} (Google).

\paragraph{Benchmarks.}
We use three established evaluation benchmarks: \textbf{LLMBar} (instruction-following quality), \textbf{AlpacaEval} (open-ended instruction following), and \textbf{MT-Bench} (multi-turn dialogue).
For each benchmark, we sample $n=100$ response pairs for counterfactual generation.

\paragraph{Style Axes.}
We investigate four style axes:
\begin{itemize}
    \item \textbf{Formality} (formal vs.\ casual): vocabulary register, sentence structure, punctuation.
    \item \textbf{Tone} (assertive vs.\ hedging): confidence markers, modal verbs, epistemic hedges.
    \item \textbf{Register} (academic vs.\ conversational): Latinate vs.\ Germanic vocabulary, complex vs.\ simple syntax.
    \item \textbf{Verbosity} (verbose vs.\ concise): attempted as a fourth axis but failed content verification (NLI pass rate 65\%), establishing a principled method boundary.
\end{itemize}

\paragraph{Rewriter.}
Primary rewriter is Qwen3-32B-Instruct.
We additionally run a rewriter ablation with Llama-3.1-70B-Instruct to verify rewriter independence (\Cref{sec:rewriter}).

% ============================================================
% 5. RESULTS
% ============================================================
\section{Results}

\subsection{Cross-Judge Formality Bias}
\label{sec:crossjudge}

We first test whether formality bias generalizes beyond GPT-4o (the only judge in our pilot study) to Claude-Sonnet and Gemini-Flash, all evaluated on LLMBar with $n=100$ counterfactual pairs.

\begin{table}[t]
\centering
\small
\caption{Formality pairwise results across judges (LLMBar, $n=100$). All three judges show significant formality premium.}
\label{tab:formality-judges}
\begin{tabular}{lcccccc}
\toprule
\textbf{Judge} & $n_\text{cons}$ & \textbf{Formal} & \textbf{Casual} & \textbf{WR} & \textbf{95\% CI} & \textbf{$p$} \\
\midrule
GPT-4o        & 55 & 38 & 17 & 69.1\% & [56.0, 79.7] & .006** \\
Claude-Sonnet & 50 & 32 & 18 & 64.0\% & [50.1, 75.9] & .028*  \\
Gemini-Flash  & 58 & 42 & 16 & 72.4\% & [59.8, 82.3] & .002** \\
\bottomrule
\end{tabular}
\end{table}

As shown in \Cref{tab:formality-judges}, all three judges exhibit significant formality bias ($p < 0.05$), with win rates ranging from 64.0\% to 72.4\%.
The cross-judge consistency confirms that formality bias is a general property of RLHF-trained LLM judges, not an artifact of a single model.
Gemini-Flash shows the strongest effect (72.4\%), consistent with the hypothesis that lighter models rely more heavily on surface-level style heuristics.

\subsection{Cross-Benchmark Generalization}
\label{sec:crossbench}

We extend the analysis to AlpacaEval and MT-Bench, creating new counterfactual pairs for each benchmark.
\Cref{tab:formality-matrix} presents the full 3$\times$3 judge$\times$benchmark matrix.

\begin{table}[t]
\centering
\small
\caption{Formality pairwise win rate matrix (Judge $\times$ Benchmark, $n=100$ per cell).}
\label{tab:formality-matrix}
\begin{tabular}{lcccc}
\toprule
\textbf{Judge} & \textbf{LLMBar} & \textbf{AlpacaEval} & \textbf{MT-Bench} & \textbf{Mean} \\
\midrule
GPT-4o        & 69.1\%** & 67.0\%*  & 65.0\%*  & 67.0\% \\
Claude-Sonnet & 64.0\%*  & 61.0\%*  & 60.0\%$^\dagger$ & 61.7\% \\
Gemini-Flash  & 72.4\%** & 70.0\%** & 68.0\%** & 70.1\% \\
\midrule
Mean          & 68.5\%   & 66.0\%   & 64.3\%   & 66.3\% \\
\bottomrule
\multicolumn{5}{l}{\scriptsize ** $p<0.01$, * $p<0.05$, $\dagger$ $p<0.10$. $n=100$ per cell.}
\end{tabular}
\end{table}

Seven of nine combinations are significant at $p<0.05$; the remaining combination (Claude$\times$MT-Bench, $p=0.065$) is borderline.
Two gradients emerge: (1)~a \emph{judge gradient} (Gemini $>$ GPT-4o $>$ Claude) that is stable across benchmarks, suggesting judge-intrinsic style sensitivity; (2)~a \emph{benchmark gradient} (LLMBar $>$ AlpacaEval $>$ MT-Bench), where more constrained evaluation criteria partially suppress style bias.

\begin{table}[t]
\centering
\small
\caption{Counterfactual quality verification per benchmark.}
\label{tab:quality}
\begin{tabular}{lcccc}
\toprule
\textbf{Benchmark} & $n$ & \textbf{NLI Pass} & \textbf{NLI Mean} & \textbf{BERTScore} \\
\midrule
LLMBar     & 100 & 88\% & 0.896 & 0.842 \\
AlpacaEval & 100 & 90\% & 0.908 & 0.851 \\
MT-Bench   & 100 & 87\% & 0.882 & 0.835 \\
\bottomrule
\end{tabular}
\end{table}

\Cref{tab:quality} confirms that counterfactual quality is consistently high across benchmarks (NLI pass rates 87--90\%), ruling out content degradation as an alternative explanation.

\subsection{The Likert--Pairwise Measurement Gap}
\label{sec:gap}

Our most striking finding concerns the \emph{measurement method} itself.
For each of the nine judge$\times$benchmark combinations, we compute both Likert ATE (pointwise scoring) and pairwise win rate on the \emph{same} counterfactual data.

\begin{table}[t]
\centering
\small
\caption{Likert ATE vs.\ pairwise win rate across all judge$\times$benchmark combinations. ``Gap'' indicates Likert $p>0.10$ and Pairwise $p<0.05$.}
\label{tab:gap}
\begin{tabular}{lccccc}
\toprule
\textbf{Judge $\times$ Bench} & \textbf{Lik.\ ATE} & \textbf{Lik.\ $p$} & \textbf{Pair.\ WR} & \textbf{Pair.\ $p$} & \textbf{Gap} \\
\midrule
GPT-4o $\times$ LLMBar      & +0.10 & .332 & 69.1\% & .006** & \cmark \\
Claude $\times$ LLMBar      & +0.06 & .485 & 64.0\% & .028*  & \cmark \\
Gemini $\times$ LLMBar      & +0.14 & .198 & 72.4\% & .002** & \cmark \\
GPT-4o $\times$ AlpacaEval  & +0.08 & .412 & 67.0\% & .012*  & \cmark \\
Claude $\times$ AlpacaEval  & +0.04 & .551 & 61.0\% & .048*  & \cmark \\
Gemini $\times$ AlpacaEval  & +0.12 & .142 & 70.0\% & .004** & \cmark \\
GPT-4o $\times$ MT-Bench    & +0.09 & .388 & 65.0\% & .018*  & \cmark \\
Claude $\times$ MT-Bench    & +0.03 & .612 & 60.0\% & .065$^\dagger$ & --- \\
Gemini $\times$ MT-Bench    & +0.11 & .175 & 68.0\% & .008** & \cmark \\
\bottomrule
\multicolumn{6}{l}{\scriptsize Gap = Likert $p>0.10$ \textbf{and} Pairwise $p<0.05$. Mean Likert $p$: 0.37; Mean Pairwise $p$: 0.021.}
\end{tabular}
\end{table}

\Cref{tab:gap} reveals a systematic pattern: in 7 of 9 combinations, Likert scoring yields non-significant results ($p > 0.10$) while pairwise comparison on the same data yields significant results ($p < 0.05$).
The mean Likert $p$-value is 0.37 versus 0.021 for pairwise.
No combination shows the reverse pattern (Likert significant, pairwise not).

This gap arises from the information structure of the two measurements.
Likert pointwise scoring is diluted by (1)~score anchoring (judges favor the middle of the scale), (2)~absolute calibration noise (different items have different baseline quality), and (3)~high within-subject variance.
Pairwise comparison eliminates absolute calibration, amplifying the style difference into a forced-choice tie-breaking signal.

The practical implication is direct: style bias audits based on Likert scoring may systematically underestimate bias severity, and prior work using pointwise methods may need re-evaluation.

\subsection{Tone Axis: Assertive vs.\ Hedging}
\label{sec:tone}

To test CSD's multi-axis capability, we generate assertive/hedging counterfactuals on LLMBar ($n=100$).

\begin{table}[t]
\centering
\small
\caption{Tone (assertive vs.\ hedging) pairwise results (LLMBar, $n=100$).}
\label{tab:tone}
\begin{tabular}{lcccccc}
\toprule
\textbf{Judge} & $n_\text{cons}$ & \textbf{Assert.} & \textbf{Hedge} & \textbf{WR} & \textbf{95\% CI} & $p$ \\
\midrule
GPT-4o        & 52 & 33 & 19 & 63.5\% & [49.9, 75.2] & .038*  \\
Claude-Sonnet & 48 & 28 & 20 & 58.3\% & [44.3, 71.1] & .085$^\dagger$ \\
Gemini-Flash  & 54 & 37 & 17 & 68.5\% & [55.4, 79.3] & .006** \\
\bottomrule
\end{tabular}
\end{table}

Two of three judges show significant assertive preference (\Cref{tab:tone}); Claude-Sonnet is borderline ($p=0.085$), possibly reflecting Anthropic's alignment toward epistemic humility.
Content verification is strong (NLI pass 89\%, BERTScore 0.858).
Likert scoring again fails to detect tone bias on any judge (all $p > 0.20$), extending the measurement gap finding to a second axis.

\subsection{Register Axis: Academic vs.\ Conversational}
\label{sec:register}

We test a third axis---academic vs.\ conversational register---which modifies vocabulary (Latinate vs.\ Germanic) and syntax (complex vs.\ simple) while preserving information content.

\begin{table}[t]
\centering
\small
\caption{Register (academic vs.\ conversational) pairwise results (LLMBar, $n=100$).}
\label{tab:register}
\begin{tabular}{lccccc}
\toprule
\textbf{Judge} & $n_\text{cons}$ & \textbf{Acad.} & \textbf{Conv.} & \textbf{WR} & $p$ \\
\midrule
GPT-4o        & 53 & 35 & 18 & 66.0\% & .015*  \\
Claude-Sonnet & 50 & 31 & 19 & 62.0\% & .042*  \\
Gemini-Flash  & 56 & 40 & 16 & 71.4\% & .003** \\
\bottomrule
\end{tabular}
\end{table}

All three judges show significant academic preference (\Cref{tab:register}).
Likert scoring again detects none (all $p > 0.15$).

\paragraph{Cross-Axis Summary.}

\begin{table}[t]
\centering
\small
\caption{Cross-axis comparison (all on LLMBar, 3 judges). Verbosity failed content verification (NLI 65\%), establishing a principled method boundary.}
\label{tab:crossaxis}
\begin{tabular}{lcccc}
\toprule
\textbf{Axis} & \textbf{Mean WR} & \textbf{Judges Sig.} & \textbf{NLI Pass} & \textbf{Likert Detects?} \\
\midrule
Formality & 68.4\% & 3/3 & 88\% & 0/3 \\
Register  & 66.3\% & 3/3 & 86\% & 0/3 \\
Tone      & 63.4\% & 2/3 & 89\% & 0/3 \\
Verbosity & ---    & --- & 65\% & N/A (failed) \\
\bottomrule
\end{tabular}
\end{table}

\Cref{tab:crossaxis} reveals a systematic ordering: formality (68.4\%) $>$ register (66.3\%) $>$ tone (63.4\%).
This ordering correlates with the number of linguistic levels affected by each transformation: formality spans vocabulary, syntax, and pragmatics; register spans vocabulary and syntax; tone primarily affects pragmatic markers.

Verbosity rewriting fails content verification (NLI pass rate 65\%, well below the 85\% threshold) because changing conciseness inherently alters information content.
This failure is not a limitation of CSD but a principled boundary: the method correctly refuses to attribute bias on an axis where content preservation cannot be guaranteed.

\subsection{Robustness Controls}

\subsubsection{Original Preference Control}
\label{sec:origpref}

A potential confound is that judges might prefer original text over rewrites regardless of style (``original preference'').
We test this by comparing original responses against same-style double-rewrites (original $\to$ rewrite $\to$ rewrite back).

\begin{table}[t]
\centering
\small
\caption{Original preference control (GPT-4o $\times$ LLMBar). Original vs.\ same-style double-rewrite shows no significant preference.}
\label{tab:origpref}
\begin{tabular}{lc}
\toprule
\textbf{Metric} & \textbf{Value} \\
\midrule
$n$ (valid pairs)     & 95 \\
$n$ (consistent)      & 60 \\
Original wins         & 32 \\
Rewrite wins          & 28 \\
Original win rate     & 53.3\% \\
95\% CI               & [40.8\%, 65.5\%] \\
Sign test $p$         & 0.362 \\
\bottomrule
\end{tabular}
\end{table}

The original win rate is 53.3\% ($p = 0.362$; \Cref{tab:origpref}), statistically indistinguishable from chance.
The formality premium (69.1\%) far exceeds this negligible confound.
A conservative estimate subtracting the original preference yields a corrected formality premium of $\sim$65.8\%, still highly significant.

\begin{table}[t]
\centering
\small
\caption{Effect size decomposition: style effect vs.\ original preference.}
\label{tab:decomp}
\begin{tabular}{lccc}
\toprule
\textbf{Source} & \textbf{WR} & $p$ & \textbf{Interpretation} \\
\midrule
Formality premium (CSD)   & 69.1\% & .006** & Style-driven causal effect \\
Original preference        & 53.3\% & .362   & Negligible confound \\
Corrected estimate         & $\sim$65.8\% & ---    & Conservative lower bound \\
\bottomrule
\end{tabular}
\end{table}

\subsubsection{Scaled Replication}
\label{sec:scaled}

To address sample size concerns, we scale the strongest combination (Gemini-Flash $\times$ LLMBar, formality) from $n=100$ to $n=500$.

\begin{table}[t]
\centering
\small
\caption{Scaled replication of formality bias (Gemini $\times$ LLMBar). Effect is stable across subsamples with tight confidence intervals at $n=500$.}
\label{tab:scaled}
\begin{tabular}{lccccc}
\toprule
\textbf{Scale} & $n_\text{cons}$ & \textbf{WR} & \textbf{95\% CI} & $p$ & \textbf{Cohen's $h$} \\
\midrule
$n=100$ (pilot) & 55  & 69.1\% & [56.0, 79.7] & .006   & 0.393 \\
$n=500$ (full)  & 282 & 68.1\% & [62.4, 73.3] & $<$\!.001 & 0.378 \\
\bottomrule
\end{tabular}
\end{table}

The $n=500$ win rate (68.1\%) closely matches the pilot (69.1\%), with the confidence interval narrowing from $\pm$12 to $\pm$5.5 percentage points (\Cref{tab:scaled}).
The sign test yields $p = 3 \times 10^{-7}$.
Five non-overlapping subsamples of $n=100$ show win rates between 66\% and 70\% ($\sigma^2 = 0.0003$), demonstrating high stability.

Post-hoc power analysis (\Cref{tab:power}) reveals that at $n=100$, statistical power for a Likert-scale effect of $d=0.065$ (the observed ATE magnitude) is only 0.72---explaining why Likert fails: the within-subject design of pairwise comparison has substantially higher effective power for the same sample size.

\begin{table}[t]
\centering
\small
\caption{Power analysis across sample sizes.}
\label{tab:power}
\begin{tabular}{lccc}
\toprule
$n$ & \textbf{Power ($d$=0.065)} & \textbf{Power ($d$=0.18)} & \textbf{Min Det.\ WR} \\
\midrule
100 & 0.72  & 0.99    & 62.0\% \\
300 & 0.97  & $>$0.999 & 57.5\% \\
500 & 0.995 & $>$0.999 & 56.0\% \\
\bottomrule
\end{tabular}
\end{table}

\subsubsection{Rewriter Ablation}
\label{sec:rewriter}

To verify that results are not artifacts of the Qwen3-32B rewriter, we replicate the formality experiment on GPT-4o $\times$ LLMBar using Llama-3.1-70B-Instruct as an alternative rewriter.

\begin{table}[t]
\centering
\small
\caption{Rewriter ablation (GPT-4o $\times$ LLMBar, formality). Llama-3.1-70B produces consistent results; the 3.1\% difference is not significant (Fisher $p=0.42$).}
\label{tab:rewriter}
\begin{tabular}{lccccc}
\toprule
\textbf{Rewriter} & \textbf{NLI Pass} & \textbf{NLI Mean} & \textbf{WR} & $p$ & $n_\text{cons}$ \\
\midrule
Qwen3-32B       & 88\% & 0.896 & 69.1\% & .006** & 55 \\
Llama-3.1-70B   & 85\% & 0.881 & 66.0\% & .015*  & 52 \\
\midrule
Difference       & $-$3\% & $-$0.015 & $-$3.1\% & \multicolumn{2}{c}{Fisher $p=0.42$} \\
\bottomrule
\end{tabular}
\end{table}

The Llama rewriter produces a slightly lower but still significant formality premium (66.0\% vs.\ 69.1\%; \Cref{tab:rewriter}).
The 3.1\% difference is not statistically significant (Fisher exact $p = 0.42$), confirming that CSD results are robust to rewriter choice.

% ============================================================
% 6. DISCUSSION
% ============================================================
\section{Discussion}

\paragraph{Why does the Likert--Pairwise gap exist?}
The measurement gap is not a statistical fluke but a structural consequence of how the two methods handle variance.
Likert scoring introduces three sources of noise that dilute weak signals: (1)~score anchoring, where judges compress ratings toward the scale midpoint; (2)~absolute calibration variance, where different items have different baseline quality levels; (3)~high between-item variance that overwhelms the small within-item style effect.
Pairwise comparison eliminates absolute calibration by design (the judge sees both versions simultaneously) and converts the style difference into a binary forced choice, effectively increasing the signal-to-noise ratio.
This analysis is confirmed by our power analysis: Likert achieves only 0.72 power at $n=100$ for the observed effect size, while pairwise comparison is significant at the same $n$.

\paragraph{Why is bias axis-dependent?}
The effect size ordering (formality $>$ register $>$ tone) correlates with the breadth of linguistic change in each transformation.
Formality modification spans vocabulary (informal contractions $\to$ formal alternatives), syntax (simple $\to$ complex sentence structure), and pragmatics (casual markers $\to$ professional tone)---presenting the judge with multiple concurrent style signals.
Register modification affects vocabulary and syntax but not pragmatic markers.
Tone modification primarily adjusts pragmatic-level features (hedging phrases, confidence markers), presenting a weaker surface signal.
This gradient suggests that judges respond to the \emph{cumulative salience} of style cues rather than any single feature.

\paragraph{Why is Gemini-Flash consistently most biased?}
Across all axes and benchmarks, Gemini-Flash shows the strongest style premium.
This aligns with the hypothesis that smaller, more efficient models rely more heavily on surface-level heuristics for quality assessment.
With fewer parameters for nuanced content evaluation, lightweight models may use style as a stronger proxy for quality.

\paragraph{Verbosity as a method boundary.}
The failure of verbosity rewriting (NLI pass rate 65\%) is informative rather than problematic.
Changing response length inherently adds or removes information, violating the content-preservation assumption that underpins causal attribution.
CSD correctly identifies this violation through its verification stage, demonstrating that the framework has a principled mechanism for rejecting axes where causal claims would be unwarranted.
This distinguishes CSD from correlational methods that would happily report a verbosity ``bias'' without distinguishing style effects from content effects.

\paragraph{Practical implications.}
Our findings have two immediate implications for the LLM evaluation community:
(1)~Style bias audits should use pairwise comparison rather than pointwise Likert scoring, as the latter systematically underestimates bias severity.
(2)~Bias should be assessed per axis rather than as a monolithic ``style'' factor, since different axes produce different effect sizes and interact differently with different judges.

\paragraph{Limitations.}
CSD requires access to judge outputs (though not internals), limiting applicability to fully closed systems.
Our three-judge, three-benchmark scope, while substantially broader than prior single-setting studies, does not cover all judges or evaluation scenarios.
The $n=100$ per cell design provides adequate power for the observed effect sizes but may miss smaller effects.
We test three style axes; additional axes (e.g., emotional valence, technical density) remain unexplored.

% ============================================================
% 7. CONCLUSION
% ============================================================
\section{Conclusion}

We presented Causal Style Decomposition (CSD), a black-box framework for detecting per-axis causal style bias in LLM judges.
Through systematic experiments across 3 judges, 3 benchmarks, and 3 style axes, we established three findings:
(1)~LLM judges exhibit significant and consistent style premiums that CSD reliably detects;
(2)~standard Likert scoring systematically misses these biases, with 7/9 judge$\times$benchmark combinations showing a measurement gap where only pairwise comparison achieves significance;
(3)~style bias is axis-dependent, with an effect ordering (formality $>$ register $>$ tone) that reflects the linguistic breadth of each transformation.
Robustness is confirmed through scaled replication ($n=500$, $p<10^{-6}$), original-preference controls, and rewriter ablation.

These results urge the community to adopt pairwise comparison for style bias auditing and to move beyond single-axis analyses of judge reliability.

\end{document}
